\documentclass[../libro.tex]{subfiles}

\begin{document}

\ifSubfilesClassLoaded{\mainmatter\chapter{行列式}\clearpage}{}

\KunAsteriskoEnEnhavtabelo
\section{\texorpdfstring{由 \(n\)~个 \(n\)~元
      \({\leq} 1\)~次方程作成的方程组 (2)}%
  {由 n 个 n 元 ≤1 次方程作成的方程组 (2)}}
\SenAsteriskoEnEnhavtabelo

\maldevigalegajxo

前面, 我们定量地研究了%
由 \(n\)~个 \(n\)~元 \({\leq} 1\)~次方程作成的方程组%
的解.
具体地, 当 \(n\)~级阵 \(A\) 的行列式不是 \(0\) 时,
我们用行列式写出了 \(AX = B\) 的唯一的解,
其中,
\(B\) 是 \(n \times 1\)~阵,
且 \(X\) 是未知的 \(n \times 1\)~阵.
不过, 若 \(\det {(A)} = 0\),
则 Cramer 公式不可用
(因为分母不可为 \(0\)).
其实, 当 \(\det {(A)} = 0\) 时,
\(AX = B\) 是否有解是一个复杂的问题.

\begin{example}
    设
    \begin{align*}
        A = \begin{bmatrix}
                1  & 1  \\
                -1 & -1 \\
            \end{bmatrix},
        \quad
        B =
        \begin{bmatrix}
            1 \\
            b \\
        \end{bmatrix},
        \quad
        X =
        \begin{bmatrix}
            x_1 \\
            x_2 \\
        \end{bmatrix},
    \end{align*}
    其中, \(b\) 是常数.
    不难算出, \(\det {(A)} = 0\).
    当 \(b\) 取某些数时,
    我们讨论 \(AX = B\) 的解.

    当 \(b = -1\) 时, 此方程组有解:
    \begin{align*}
        A
        \begin{bmatrix}
            1 \\
            0 \\
        \end{bmatrix}
        =
        \begin{bmatrix}
            1 \cdot 1 + 1 \cdot 0     \\
            -1 \cdot 1 + (-1) \cdot 0 \\
        \end{bmatrix}
        =
        B.
    \end{align*}
    进一步地, \(AX = B\) 的解不唯一.
    % The following (commented) text
    % assumes that the underlying field
    % has infinitely many elements.
    % 设 \(t\) 是常数.
    % 则
    % \begin{align*}
    %     A
    %     \begin{bmatrix}
    %         1 - t \\
    %         t     \\
    %     \end{bmatrix}
    %     =
    %     \begin{bmatrix}
    %         1 \cdot (1 - t) + 1 \cdot t     \\
    %         -1 \cdot (1 - t) + (-1) \cdot t \\
    %     \end{bmatrix}
    %     =
    %     B.
    % \end{align*}
    % 取不同的 \(t\),
    % 就有不一样的
    % \(
    % \begin{bmatrix}
    %     1 - t \\
    %     t     \\
    % \end{bmatrix}
    % \).
    % 我们知道, 有无限多个数.
    % 所以, \(AX = B\) 有无限多个解.
    不难算出,
    \begin{align*}
        A
        \begin{bmatrix}
            0 \\
            1 \\
        \end{bmatrix}
        =
        \begin{bmatrix}
            1 \cdot 0 + 1 \cdot 1     \\
            -1 \cdot 0 + (-1) \cdot 1 \\
        \end{bmatrix}
        =
        B,
    \end{align*}
    且
    \(
    \begin{bmatrix}
        0 \\
        1 \\
    \end{bmatrix}
    \neq
    \begin{bmatrix}
        1 \\
        0 \\
    \end{bmatrix}
    \).

    当 \(b = 0\) 时, 此方程组无解.
    用反证法.
    反设
    \(
    A \begin{bmatrix}
        y_1 \\
        y_2 \\
    \end{bmatrix}
    = B
    \),
    即
    \begin{align*}
        y_1 + y_2  & = 1, \\
        -y_1 - y_2 & = 0.
    \end{align*}
    从而
    \begin{align*}
        1 + 0
        = {} & (y_1 + y_2) + (-y_1 - y_2)     \\
        = {} & (y_1 - y_1) + (y_2 - y_2) = 0.
    \end{align*}
    这是矛盾.
\end{example}

从本节开始, 我们定性地研究%
由 \(n\)~个 \(n\)~元 \({\leq} 1\)~次方程作成的方程组%
的解.
具体地, 我们主要讨论解的性质,
而不是解的公式.

根据 Cramer 公式, 我们不难得到如下事实:

\begin{theorem}
    设 \(A\) 是 \(n\)~级阵.
    设 \(\det {(A)} \neq 0\).
    那么, 对任何的 \(n \times 1\)~阵 \(B\),
    存在一个 \(n \times 1\)~阵 \(C\),
    使 \(AC = B\).
    (或者,
    若存在某 \(n \times 1\)~阵 \(B\),
    使对任何的 \(n \times 1\)~阵 \(C\),
    必 \(AC \neq B\),
    则 \(\det {(A)} = 0\).)
\end{theorem}

重要地, 此事反过来也对.

\begin{theorem}
    设 \(A\) 是 \(n\)~级阵.
    设对任何的 \(n \times 1\)~阵 \(B\),
    存在一个 \(n \times 1\)~阵 \(C\),
    使 \(AC = B\).
    则 \(\det {(A)} \neq 0\).
    (或者,
    若 \(\det {(A)} = 0\),
    则存在某 \(n \times 1\)~阵 \(B\),
    使对任何的 \(n \times 1\)~阵 \(C\),
    必 \(AC \neq B\).)
\end{theorem}

\begin{proof}
    设 \(n\)~级单位阵~\(I\) 的%
    列~\(1\), \(2\), \(\dots\), \(n\)
    分别是
    \(e_1\), \(e_2\), \(\dots\), \(e_n\).
    每一个 \(e_i\) 都是 \(n \times 1\)~阵.
    根据假定, 存在一个 \(n \times 1\)~阵 \(f_i\)
    使 \(Af_i = e_i\).
    作 \(n\)~级阵
    \(F = [f_1, f_2, \dots, f_n]\).
    则
    \begin{align*}
        AF
        = {} & A\, [f_1, f_2, \dots, f_n] \\
        = {} & [Af_1, Af_2, \dots, Af_n]  \\
        = {} & [e_1, e_2, \dots, e_n]     \\
        = {} & I.
    \end{align*}
    因为二个同级的方阵的积的行列式%
    等于这二个阵的行列式的积,
    \begin{align*}
        1 = \det {(I)} = \det {(A)} \det {(F)}.
    \end{align*}
    从而 \(\det {(A)} \neq 0\).
\end{proof}

\KunAsteriskoEnEnhavtabelo
\section{\texorpdfstring{由 \(n\)~个 \(n\)~元
      \({\leq} 1\)~次方程作成的方程组 (3)}%
  {由 n 个 n 元 ≤1 次方程作成的方程组 (3)}}
\SenAsteriskoEnEnhavtabelo

\maldevigalegajxo

前面, 我们知道,
若 \(n\)~级阵 \(A\) 的行列式为 \(0\),
则存在某 \(n \times 1\)~阵 \(B\),
使线性方程组 \(AX = B\) 无解.
当然, 对某些 \(B\),
\(AX = B\) 还是有解的:
取 \(B = 0\),
则 \(AX = B\) 至少有一个解 \(X = 0\)
(\(0\) 是元全为 \(0\) 的 \(n \times 1\)~阵).
本节, 我们研究, 若 \(AX = B\) 有解,
则它是否有唯一的解.

我们先看一个简单的事实.

\begin{theorem}
    设 \(A\) 是 \(n\)~级阵.
    设 \(B\) 是 \(n \times 1\)~阵.
    设 \(n \times 1\)~阵 \(C\) 适合 \(AC = B\).

    (1)
    若 \(AX = 0\) 只有零解, 则 \(AX = B\) 的解唯一;

    (2)
    若存在非零的 \(n \times 1\)~阵 \(D\)
    使 \(AD = 0\),
    则 \(AX = B\) 的解不唯一.
\end{theorem}

\begin{proof}
    (1)
    设 \(n \times 1\)~阵 \(Y\) 适合 \(AY = B\).
    则
    \begin{align*}
        A(Y - C) = AY - AC = B - B = 0.
    \end{align*}
    故 \(Y - C\) 是 \(AX = 0\) 的一个解.
    因为 \(AX = 0\) 只有零解,
    故 \(Y - C = 0\).
    从而 \(Y = C\).

    % The following proof assumes that the underlying field
    % has infinitely many elements.

    % (2)
    % 要证 \(AX = B\) 有无限多个解,
    % 只要证,
    % 对任何正整数 \(m\),
    % 我们能找到 \(AX = B\) 的
    % \(m\) 个互不相同的解.

    % 作 \(m\)~个 \(n \times 1\)~阵
    % \(C_1\), \(C_2\), \(\dots\), \(C_m\),
    % 其中, \(C_j = C + jD\).
    % 首先, 每一个 \(C_j\) 都是 \(AX = B\) 的解:
    % \begin{align*}
    %     AC_j = A(C + jD) = AC + A(jD) = B + j(AD) = B + j0 = B.
    % \end{align*}
    % 然后, 我们证,
    % 若 \(p\), \(q\) 是二个不相等的%
    % 且不超过 \(m\) 的正整数,
    % 则 \(C_p \neq C_q\).
    % 用反证法.
    % 反设 \(C_p = C_q\),
    % 则
    % \begin{align*}
    %     0 = C_p - C_q = (C + pD) - (C + qD) = (p - q)D.
    % \end{align*}
    % 因为 \(p \neq q\), 故 \(p - q \neq 0\).
    % 从而
    % \begin{align*}
    %     0 = \frac{1}{p-q}\, 0
    %     = \frac{1}{p-q}\, ((p-q)D)
    %     = \Big( \frac{1}{p-q}\, (p-q) \Big) D
    %     = 1D
    %     = D.
    % \end{align*}
    % 这是矛盾.
    % 所以, \(C_p \neq C_q\).

    (2)
    因为 \(D \neq 0\), 故 \(C + D \neq C\).
    因为 \(A(C + D) = AC + AD = B + 0 = B\),
    且 \(C + D \neq C\),
    故 \(AX = B\) 的解不唯一.
\end{proof}

由此可见, 若 \(AX = 0\) 有非零解,
则 \(AX = B\) 有解时,
它的解不唯一.
若 \(AX = 0\) 只有零解,
则 \(AX = B\) 有解时,
它的解唯一.

根据 Cramer 公式, 我们不难得到如下事实:

\begin{theorem}
    设 \(A\) 是 \(n\)~级阵.
    设 \(AX = 0\) 有非零解.
    则 \(\det {(A)} = 0\).
    (或者,
    若 \(\det {(A)} \neq 0\),
    则 \(AX = 0\) 只有零解.)
\end{theorem}

重要地, 此事反过来也对.

\begin{theorem}
    设 \(A\) 是 \(n\)~级阵.
    设 \(\det {(A)} = 0\).
    则 \(AX = 0\) 有非零解.
    (或者,
    若 \(AX = 0\) 只有零解,
    则 \(\det {(A)} \neq 0\).)
\end{theorem}

为方便, 我要引入一个小概念.
说阵~\(Z\) 是 \(A\) 的一个 \emph{\(p\)~级子阵},
就是说, \(Z\) 是 \(A\) 的一个子阵
(见本章, 节~\sekcio{5}),
且 \(Z\) 是一个 \(p\)~级阵.

在论证此事前, 我想用三个例%
助您理解此事为什么是对的.

\begin{example}
    设 \(A\) 是一个 \(5\)~级阵,
    且 \(A = 0\).
    那么, 对任何非零的 \(5 \times 1\)~阵 \(S\),
    必有 \(AS = 0\).
\end{example}

\begin{example}
    设 \(A\) 是一个 \(5\)~级阵,
    且 \(\det {(A)} = 0\),
    但 \(A \neq 0\).

    我们知道,
    \(A \operatorname{adj} {(A)} = \det {(A)}\, I\),
    其中, \(\operatorname{adj} {(A)}\)
    是 \(A\) 的古伴,
    且 \(I\) 是 \(5\)~级单位阵.
    因为 \(\det {(A)} = 0\),
    故 \(A \operatorname{adj} {(A)} = 0\).

    若 \(\operatorname{adj} {(A)} \neq 0\),
    则存在 \(k\), \(v\), 使
    \([\operatorname{adj} {(A)}]_{k,v} \neq 0\).
    我们设 \(\operatorname{adj} {(A)}\) 的%
    列~\(v\) 是 \(Y\).
    那么, \([Y]_{k,1}
        = [\operatorname{adj} {(A)}]_{k,v} \neq 0\),
    从而 \(Y \neq 0\).
    我们证明 \(AY = 0\):
    \begin{align*}
        [AY]_{i,1}
        = {} &
        \sum_{\ell = 1}^{5}
        {[A]_{i,\ell} [Y]_{\ell,1}}
        \\
        = {} &
        \sum_{\ell = 1}^{5}
        {[A]_{i,\ell} [\operatorname{adj} {(A)}]_{\ell,v}}
        \\
        = {} &
        [A \operatorname{adj} {(A)}]_{i,v}
        \\
        = {} &
        [0]_{i,v}
        \\
        = {} &
        0.
    \end{align*}
\end{example}

\begin{example}
    仍设 \(A\) 是一个 \(5\)~级阵,
    且 \(\det {(A)} = 0\),
    但 \(A \neq 0\).

    在上个例里, 我们假定
    \(\operatorname{adj} {(A)} \neq 0\).
    可是, 它也有可能为 \(0\), 即使 \(A \neq 0\)
    (比如, 可以验证
    \begin{align*}
        M = \begin{bmatrix}
                1  & 1 & -1 & -1 & 1 \\
                1  & 2 & 0  & 0  & 2 \\
                2  & 1 & -3 & -3 & 1 \\
                -1 & 0 & 2  & 2  & 0 \\
                -1 & 1 & 3  & 3  & 1 \\
            \end{bmatrix}
    \end{align*}
    的古伴就是 \(0\)).

    现在, 我们假定
    \(\operatorname{adj} {(A)} = 0\).
    此时, 我们不能利用
    \(\operatorname{adj} {(A)}\)
    写出 \(AX = 0\) 的非零解.
    我们应如何作?

    若 \(\operatorname{adj} {(A)} = 0\),
    则 \(A\) 的每一个 \(4\)~级子阵的行列式都是 \(0\).
    从而, 存在一个%
    低于 \(4\) 的正整数 \(r\)
    适合性质:
    \(A\) 有一个 \(r\)~级子阵, 其的行列式非零,
    且 \(A\) 的每一个 \(r+1\)~级子阵的行列式都是 \(0\).
    (我们已知 \(A\) 的每一个 \(4\)~级子阵的行列式都是 \(0\).
    我们考虑 \(A\) 的 \(3\)~级子阵.
    若有一个 \(3\)~级子阵的行列式不是 \(0\),
    我们就取 \(r = 3\).
    若不然, \(A\) 的每一个 \(3\)~级子阵的行列式都是 \(0\).
    我们考虑 \(A\) 的 \(2\)~级子阵.
    若有一个 \(2\)~级子阵的行列式不是 \(0\),
    我们就取 \(r = 2\).
    若不然, \(A\) 的每一个 \(2\)~级子阵的行列式都是 \(0\).
    我们考虑 \(A\) 的 \(1\)~级子阵.
    因为 \(A \neq 0\),
    故 \(A\) 有 \(1\)~级子阵的行列式不是 \(0\).
    此时, 取 \(r = 1\) 即可.)

    我以 \(r = 2\) 为例%
    演示找 \(AX = 0\) 的非零解的方法
    (可以验证,
    前面的 \(M\) 的行~\(1\), \(2\)
    与列~\(1\), \(2\) 作成的
    \(2\)~级子阵的行列式不是 \(0\),
    但 \(M\) 的每一个 \(3\)~级子阵的行列式都是 \(0\));
    \(r = 3\) 或 \(r = 1\) 时的情形是类似的.

    我们设
    \(\displaystyle
    A_2 =
    A\binom{1,2}{1,2}
    \)
    的行列式不是 \(0\),
    且 \(A\) 的每一个 \(3\)~级子阵的行列式都是 \(0\).
    再设
    \(\displaystyle
    E =
    A\binom{1,2,3}{1,2,3}
    \).
    因为 \(E\) 是 \(A\) 的一个 \(3\)~级子阵,
    故 \(\det {(E)} = 0\).
    从而
    \(E \operatorname{adj} {(E)} = 0\).
    注意,
    \([\operatorname{adj} {(E)}]_{3,3}
    = (-1)^{3+3} \det {(A_2)} \neq 0\).
    设 \(\operatorname{adj} {(E)}\)
    的列~\(3\) 为 \(T\).
    那么 \(T \neq 0\).
    并且, 我们可用完全类似的方法, 证明
    \(ET = 0\).

    接下来, 我们设法由此得到一个 \(AX = 0\)
    的非零解.
    我们作一个 \(5 \times 1\)~阵 \(W\), 其中,
    \begin{align*}
        [W]_{\ell,1}
        = \begin{cases}
              [\operatorname{adj} {(E)}]_{\ell,3},
                 & \ell \leq 3; \\
              0, & \ell > 3.
          \end{cases}
    \end{align*}
    (通俗地, 我们在 \(T\) 的最后一个元后加 \(2\)~个 \(0\),
    变 \(3 \times 1\)~阵 \(T\)
    为一个 \(5 \times 1\)~阵 \(W\).)
    因为 \([W]_{3,1} \neq 0\),
    故 \(W \neq 0\).
    我们证明 \(AW = 0\).

    取不超过 \(5\) 的正整数 \(q\).
    则
    \begin{align*}
        [AW]_{q,1}
        = {} &
        \sum_{\ell = 1}^{5}
        {[A]_{q,\ell} [W]_{\ell,1}}
        \\
        = {} &
        \sum_{\ell = 1}^{3}
        {[A]_{q,\ell} [W]_{\ell,1}}
        +
        \sum_{\ell = 4}^{5}
        {[A]_{q,\ell} [W]_{\ell,1}}
        \\
        = {} &
        \sum_{\ell = 1}^{3}
        {[A]_{q,\ell} [\operatorname{adj} {(E)}]_{\ell,3}}
        +
        \sum_{\ell = 4}^{5}
        {[A]_{q,\ell}\, 0}
        \\
        = {} &
        \sum_{\ell = 1}^{3}
        {[A]_{q,\ell} [\operatorname{adj} {(E)}]_{\ell,3}}.
    \end{align*}
    若 \(q \leq 3\), 则
    \begin{align*}
        [AW]_{q,1}
        = {} &
        \sum_{\ell=1}^{3}
        {[A]_{q,\ell}
            [\operatorname{adj} {(E)}]_{\ell,3}}
        \\
        = {} &
        \sum_{p=1}^{3}
        {[E]_{q,\ell}
            [\operatorname{adj} {(E)}]_{\ell,3}}
        \\
        = {} &
        [E \operatorname{adj} {(E)}]_{q,3}
        \\
        = {} &
        [0]_{q,3}
        \\
        = {} & 0.
    \end{align*}
    设 \(q > 3\).
    则
    \begin{align*}
        [AW]_{q,1}
        = {} &
        \sum_{\ell=1}^{3}
        {[A]_{q,\ell}
            [\operatorname{adj} {(E)}]_{\ell,3}}
        \\
        = {} &
        \sum_{\ell=1}^{3}
        {[A]_{q,\ell}
        (-1)^{3+\ell} \det {(E(3|\ell))}}.
    \end{align*}
    作一个 \(3\)~级阵 \(C_q\), 其中,
    \begin{align*}
        [C_q]_{h,\ell}
        = \begin{cases}
              [E]_{h,\ell}, & h \neq 3; \\
              [A]_{q,\ell}, & h = 3.
          \end{cases}
    \end{align*}
    于是, \(C_q(3|\ell) = E(3|\ell)\).
    从而
    \begin{align*}
        [AW]_{q,1}
        = {} &
        \sum_{\ell=1}^{3}
        {[A]_{q,\ell}
        (-1)^{3+\ell} \det {(E(3|\ell))}}
        \\
        = {} &
        \sum_{\ell=1}^{3}
        {[C_q]_{3,\ell}
        (-1)^{3+\ell} \det {(C_q (3|\ell))}}
        \\
        = {} &
        \det {(C_q)}.
    \end{align*}
    注意,
    \(
    \displaystyle
    C_q = A\binom{1,2,q}{1,2,3}
    \)
    是 \(A\)~的一个 \(3\)~级子阵,
    故
    \begin{align*}
        [AW]_{q,1} = \det {(C_q)} = 0.
    \end{align*}

    综上, \(AW = 0\), 且 \(W \neq 0\).
\end{example}

下面, 我给出此事的论证.
证明分三个情形:
(a)
\(A = 0\);
(b)
\(A \neq 0\), 且
\(\operatorname{adj} {(A)} \neq 0\);
(c)
\(A \neq 0\), 且
\(\operatorname{adj} {(A)} = 0\).
情形 (a) 最简单:
我们随便找一个非零的 \(n \times 1\)~阵即可.
情形 (b) 不难:
我们取 \(A\) 的古伴的一个非零的列即可.
情形 (c) 是最复杂的:
在上个例里, 我们假定%
由 \(A\) 的前 \(r\)~行与前 \(r\)~列作成的
\(r\)~级子阵的行列式不是 \(0\)
(通俗地, \(A\) 在左上角有行列式非零的 \(r\)~级子阵),
且每一个 \(r+1\)~级子阵的行列式都是 \(0\)
(此处 \(r = 2\));
不过, 一般的阵不一定在左上角有行列式非零的 \(r\)~级子阵,
且每一个 \(r+1\)~级子阵的行列式都是 \(0\);
这是论证的最大的挑战.

\begin{proof}
    设 \(A\) 是 \(n\)~级阵,
    且 \(\det {(A)} = 0\).

    若 \(A = 0\),
    我们任取一个非零的 \(n \times 1\)~阵 \(S\).
    则 \(AS = 0\).

    以下, 我们假定 \(A \neq 0\);
    也就是, \(A\) 有一个元不是 \(0\)
    (同时, 这也要求 \(n > 1\):
    回想 \(1\)~级阵的行列式是什么).

    我们知道,
    \(A \operatorname{adj} {(A)} = \det {(A)}\, I\),
    其中, \(\operatorname{adj} {(A)}\)
    是 \(A\) 的古伴,
    且 \(I\) 是 \(n\)~级单位阵.
    因为 \(\det {(A)} = 0\),
    故 \(A \operatorname{adj} {(A)} = 0\).

    若 \(\operatorname{adj} {(A)} \neq 0\),
    则存在 \(k\), \(v\), 使
    \([\operatorname{adj} {(A)}]_{k,v} \neq 0\).
    我们设 \(\operatorname{adj} {(A)}\) 的%
    列~\(v\) 是 \(Y\).
    那么, \([Y]_{k,1}
        = [\operatorname{adj} {(A)}]_{k,v} \neq 0\),
    从而 \(Y \neq 0\).
    我们证明 \(AY = 0\):
    \begin{align*}
        [AY]_{i,1}
        = {} &
        \sum_{\ell = 1}^{n}
        {[A]_{i,\ell} [Y]_{\ell,1}}
        \\
        = {} &
        \sum_{\ell = 1}^{n}
        {[A]_{i,\ell} [\operatorname{adj} {(A)}]_{\ell,v}}
        \\
        = {} &
        [A \operatorname{adj} {(A)}]_{i,v}
        \\
        = {} &
        [0]_{i,v}
        \\
        = {} &
        0.
    \end{align*}

    若 \(\operatorname{adj} {(A)} = 0\),
    则 \(A\) 的每一个 \(n-1\)~级子阵的行列式都是 \(0\)
    (同时, 这也要求 \(n > 2\):
    回想 \(2\)~级阵的古伴是什么;
    再回想前面作过的假定 \(A \neq 0\)).
    从而, 存在一个%
    低于 \(n-1\) 的正整数 \(r\)
    适合性质:
    \(A\) 有一个 \(r\)~级子阵, 其的行列式非零,
    且 \(A\) 的每一个 \(r+1\)~级子阵的行列式都是 \(0\).
    (我们已知 \(A\) 的每一个 \(n-1\)~级子阵的行列式都是 \(0\).
    我们考虑 \(A\) 的 \(n-2\)~级子阵.
    若有一个 \(n-2\)~级子阵的行列式不是 \(0\),
    我们就取 \(r = n-2\).
    若不然, \(A\) 的每一个 \(n-2\)~级子阵的行列式都是 \(0\).
    我们考虑 \(A\) 的 \(n-3\)~级子阵.
    若有一个 \(n-3\)~级子阵的行列式不是 \(0\),
    我们就取 \(r = n-3\).
    ……
    若不然, \(A\) 的每一个 \(2\)~级子阵的行列式都是 \(0\).
    我们考虑 \(A\) 的 \(1\)~级子阵.
    因为 \(A \neq 0\),
    故 \(A\) 有 \(1\)~级子阵的行列式不是 \(0\).
    此时, 取 \(r = 1\) 即可.)

    我们设 \(A\)~的 \(r\)~级子阵
    \(
    \displaystyle
    A_r = A\binom{i_1,\dots,i_r}{j_1,\dots,j_r}
    \)
    的行列式非零,
    其中, \(i_1 < i_2 < \dots < i_r\),
    且 \(j_1 < j_2 < \dots < j_r\).
    我们从 \(1\), \(2\), \(\dots\), \(n\) 中%
    去除 \(r\)~个整数
    \(i_1\), \(i_2\), \(\dots\), \(i_r\);
    此时, 还剩下 \(n-r\)~个整数;
    我们从中选一个为 \(i_{r+1}\).
    类似地,
    我们再从 \(1\), \(2\), \(\dots\), \(n\) 中%
    去除 \(r\)~个整数
    \(j_1\), \(j_2\), \(\dots\), \(j_r\);
    此时, 还剩下 \(n-r\)~个整数;
    我们从中选一个为 \(j_{r+1}\).
    作 \(r+1\)~级阵
    \(
    \displaystyle
    E =
    A\binom{i_1,\dots,i_r,i_{r+1}}
    {j_1,\dots,j_r,j_{r+1}}
    \).
    我们设 \(i_p\) 在
    \(i_1\), \(\dots\), \(i_r\), \(i_{r+1}\)
    中是第~\(f(i_p)\) 小的数.
    再设 \(j_p\) 在
    \(j_1\), \(\dots\), \(j_r\), \(j_{r+1}\)
    中是第~\(g(j_p)\) 小的数.
    因为 \(E\) 是 \(A\) 的一个 \(r+1\)~级子阵,
    故 \(\det {(E)} = 0\).
    从而
    \(E \operatorname{adj} {(E)} = 0\).
    注意,
    \([\operatorname{adj} {(E)}]_{g(j_{r+1}), f(i_{r+1})}
    = (-1)^{f(i_{r+1})+g(j_{r+1})} \det {(A_r)} \neq 0\).
    设 \(\operatorname{adj} {(E)}\)
    的列~\(f(i_{r+1})\) 为 \(T\).
    那么 \(T \neq 0\).
    并且, 我们可用完全类似的方法, 证明
    \(ET = 0\).

    接下来, 我们设法由此得到一个 \(AX = 0\)
    的非零解.
    我们作一个 \(n \times 1\)~阵 \(W\), 其中,
    \begin{align*}
        [W]_{\ell,1}
        = \begin{cases}
              [\operatorname{adj} {(E)}]_{g(j_p),f(i_{r+1})},
                 & \text{\(\ell = j_p\),
              \(p = 1\), \(\dots\), \(r\), \(r+1\)}; \\
              0, & \text{其他}.
          \end{cases}
    \end{align*}
    (通俗地, 我们在适当的位置写 \(0\),
    变 \((r+1) \times 1\)~阵 \(T\)
    为一个 \(n \times 1\)~阵 \(W\).)
    因为 \([W]_{j_{r+1},1} \neq 0\),
    故 \(W \neq 0\).
    我们证明 \(AW = 0\).

    取不超过 \(n\) 的正整数 \(q\).
    则
    \begin{align*}
        [AW]_{q,1}
        = {} &
        \sum_{\ell = 1}^{n}
        {[A]_{q,\ell} [W]_{\ell,1}}
        \\
        = {} &
        \sum_{\substack{
        1 \leq \ell \leq n \\
        \ell = j_p         \\
                1 \leq p \leq r+1
            }}
        {[A]_{q,\ell} [W]_{\ell,1}}
        +
        \sum_{\substack{
        1 \leq \ell \leq n \\
                \text{其他}
            }}
        {[A]_{q,\ell} [W]_{\ell,1}}
        \\
        = {} &
        \sum_{\substack{
        1 \leq \ell \leq n \\
        \ell = j_p         \\
                1 \leq p \leq r+1
            }}
        {[A]_{q,j_p}
                [\operatorname{adj} {(E)}]_{g(j_p),f(i_{r+1})}}
        +
        \sum_{\substack{
        1 \leq \ell \leq n \\
                \text{其他}
            }}
        {[A]_{q,\ell}\, 0}
        \\
        = {} &
        \sum_{p=1}^{r+1}
        {[A]_{q,j_p}
            [\operatorname{adj} {(E)}]_{g(j_p),f(i_{r+1})}}.
    \end{align*}
    若 \(q\) 等于某个 \(i_t\), 则
    \begin{align*}
        [AW]_{i_t,1}
        = {} &
        \sum_{p=1}^{r+1}
        {[A]_{i_t,j_p}
            [\operatorname{adj} {(E)}]_{g(j_p),f(i_{r+1})}}
        \\
        = {} &
        \sum_{p=1}^{r+1}
        {[E]_{f(i_t),g(j_p)}
            [\operatorname{adj} {(E)}]_{g(j_p),f(i_{r+1})}}
        \\
        = {} &
        [E \operatorname{adj} {(E)}]_{f(i_t),f(i_{r+1})}
        \\
        = {} &
        [0]_{f(i_t),f(i_{r+1})}
        \\
        = {} & 0.
    \end{align*}
    若 \(q\) 不等于
    \(i_1\), \(\dots\), \(i_r\), \(i_{r+1}\)
    中的任何一个,
    则
    \begin{align*}
        [AW]_{q,1}
        = {} &
        \sum_{p=1}^{r+1}
        {[A]_{q,j_p}
            [\operatorname{adj} {(E)}]_{g(j_p),f(i_{r+1})}}
        \\
        = {} &
        \sum_{p=1}^{r+1}
        {[A]_{q,j_p}
        (-1)^{f(i_{r+1})+g(j_p)}
        \det {(E(f(i_{r+1})|g(j_p)))}}.
    \end{align*}
    作一个 \(r+1\)~级阵 \(C_q\), 其中,
    \begin{align*}
        [C_q]_{h,g(j_p)}
        = \begin{cases}
              [E]_{h,g(j_p)},
               & h \neq f(i_{r+1}); \\
              [A]_{q,j_p},
               & h = f(i_{r+1}).
          \end{cases}
    \end{align*}
    于是, \(C_q(f(i_{r+1})|g(j_p))
    = E(f(i_{r+1})|g(j_p))\).
    从而
    \begin{align*}
        [AW]_{q,1}
        = {} &
        \sum_{p=1}^{r+1}
        {[A]_{q,j_p}
        (-1)^{f(i_{r+1})+g(j_p)}
        \det {(E(f(i_{r+1})|g(j_p)))}}
        \\
        = {} &
        \sum_{p=1}^{r+1}
        {[C_q]_{f(i_{r+1}),g(j_p)}
        (-1)^{f(i_{r+1})+g(j_p)}
        \det {(C_q (f(i_{r+1})|g(j_p)))}}
        \\
        = {} &
        \det {(C_q)}.
    \end{align*}

    再作一个 \(r+1\)~级阵
    \(
    \displaystyle
    D_q =
    A\binom{i_1,\dots,i_r,q}
    {j_1,\dots,j_r,j_{r+1}}
    \).
    不难看出, 适当地交换 \(C_q\) 的行的次序,
    即可变 \(C_q\) 为 \(D_q\).
    根据反称性,
    \(\det {(C_q)} = \pm \det {(D_q)}\).
    (具体地, 设 \(q\) 在
    \(i_1\), \(\dots\), \(i_r\), \(q\)
    中是第~\(u\) 小的数.
    那么, 当 \(f(i_{r+1}) = u\) 时, \(C_q = D_q\).
    当 \(f(i_{r+1}) < u\) 时,
    交换行~\(f(i_{r+1})\) 与 \(f(i_{r+1})+1\),
    再交换行~\(f(i_{r+1})+1\) 与 \(f(i_{r+1})+2\),
    ……
    且再交换行~\(u-1\) 与 \(u\);
    作 \(u - f(i_{r+1})\)~次相邻行的交换,
    即可变 \(C_q\) 为 \(D_q\).
    当 \(f(i_{r+1}) > u\) 时,
    交换行~\(f(i_{r+1})\) 与 \(f(i_{r+1})-1\),
    再交换行~\(f(i_{r+1})-1\) 与 \(f(i_{r+1})-2\),
    ……
    且再交换行~\(u+1\) 与 \(u\);
    作 \(f(i_{r+1}) - u\)~次相邻行的交换,
    即可变 \(C_q\) 为 \(D_q\).)
    因为 \(D_q\) 是 \(A\) 的一个 \(r+1\)~级子阵,
    故它的行列式为 \(0\).
    从而
    \begin{align*}
        [AW]_{q,1}
        = \det {(C_q)}
        = \pm \det {(D_q)}
        = 0.
    \end{align*}

    综上, 我们找到了一个 \(n \times 1\)~阵 \(W\)
    使 \(AW = 0\), 且 \(W \neq 0\).
\end{proof}

\KunAsteriskoEnEnhavtabelo
\section{\texorpdfstring{由 \(n\)~个 \(n\)~元
      \({\leq} 1\)~次方程作成的方程组 (4)}%
  {由 n 个 n 元 ≤1 次方程作成的方程组 (4)}}
\SenAsteriskoEnEnhavtabelo

\maldevigalegajxo

综合前几节的讨论, 我们有如下定理.

\begin{theorem}
    设 \(A\) 是 \(n\)~级阵,
    \(B\) 是 \(n \times 1\)~阵,
    且 \(X\) 是未知的 \(n \times 1\)~阵.

    (1)
    当 \(\det {(A)} \neq 0\) 时,
    \(AX = B\) 有唯一的解;

    (2)
    当 \(\det {(A)} = 0\) 时,
    \(AX = B\) 要么无解,
    要么有多于一个解.
\end{theorem}

这就是%
由 \(n\)~个 \(n\)~元 \({\leq} 1\)~次方程作成的方程组%
的解的定性的理论 (青春版).
不过,
% 不过:
% (a)
当 \(\det {(A)} = 0\) 时,
如何进一步地判断 \(AX = B\) 是否有解?
这是此理论无法解答的问题.
% (b)
% 此理论能否被推广到一般的%
% 由 \(m\)~个 \(n\)~元 \({\leq} 1\)~次方程作成的方程组?
% (c)
% 若 \(AX = B\) 有解,
% 但 \(\det {(A)} = 0\),
% 能否给出解的公式?
% 这是定量的理论.

若 \(\det {(A)} \neq 0\),
则, 定量地,
我们可用 Cramer 公式写出
\(AX = B\) 的 (唯一的) 解.
不过, 若 \(AX = B\),
但 \(\det {(A)} = 0\),
我们如何写出它的 (所有的) 解?

我会在接下来的几节里,
讨论一般的%
由 \(m\)~个 \(n\)~元 \({\leq} 1\)~次方程作成的方程组.
这些问题会被好地解决.

\end{document}

This material discusses
systems of n linear equations with n unknowns.
This material should have been deprecated,
but for some reason, I still keep it here.
All sections in this material are optional,
so it does not really matter.

The most interesting section is perhaps
the first one in this file:
it discusses uniqueness of solutions.
The other sections are boring,
because I will give generalisations of the results
in those sections
elsewhere
(or rather, in the material which discusses
systems of m linear equations with n unknowns).
Here is my lesson:
to study systems of n linear equations with n unknowns,
do not just study systems of n linear equations with n unknowns;
study systems of m linear equations with n unknowns,
and it will be clearer.
